% Proposed B.118. Owned standalone fragment; no central assembly/register edit.
\subsection*{The actual full boundary and its neck intersection}

\statement{B.118}{Proved}{The actual regional boundary subgroup is a fully embedded profinite free product, with an exact discrete and completed neck intersection}
\label{stmt:B118}
Use the actual annular regional marking of B.112--B.114, before the
two additional cross relations. Write
\[
 B=CbC,\qquad z=Ch,\qquad f=hk^{-1}=Czk^{-1},
 \qquad K=F(a,B,d,e,f)\leq R.
\]
Let $c=aCa^{-1}$ be the selected upper-gap cone, and put
\begin{equation}\label{eq:B118-marked-boundary}
 P=P_5=O_1O_2O_3O_4O_5=aBfde,\qquad
 T=cP=abzde k^{-1},\qquad H=\langle c,k,T\rangle\leq R.
\end{equation}
Then the indicated generators give an isomorphism
\begin{equation}\label{eq:B118-H}
 H\cong C_2(c)*\mathbb Z^2(k,T),
\end{equation}
and $H$ has its full induced profinite topology in $R$.
The actual neck intersection is
\begin{equation}\label{eq:B118-intersections}
 K\cap H=\langle P\rangle,\qquad
 \widehat K\cap\widehat H
       =\overline{\langle P\rangle}\quad\text{inside }\widehat R.
\end{equation}
Here the hats are full profinite completions embedded by the proved
induced-topology comparisons. The bar denotes the actual compact
closure in $\widehat R$; it is the full procyclic completion of
$\langle P\rangle$.

There is also an auxiliary corner group
\begin{equation}\label{eq:B118-corner}
 W_{\partial}=
 \langle D,u,k,j\mid D^2=j^2=1,\ [k,D]=[k,u]=1,
                         \ juj=u^{-1}\rangle.
\end{equation}
The marked assignment $c\mapsto jDj$, $T\mapsto uD$, $k\mapsto k$
embeds the free product \eqref{eq:B118-H} in $W_{\partial}$, again
with its full induced profinite topology. This is a comparison with
the isolated corner: no embedding after the remaining regional or
mirror relations, and no properness assertion for an additional
profinite amalgam, is included in the theorem.

\begin{proof}
\emph{B.118.1. The actual boundary and a retaining quotient.}
The physical neck basis of B.114 has
\[
 (O_1,O_2,O_3,O_4,O_5)
 =(aBa^{-1},a,fde d^{-1}f^{-1},d,d^{-1}fd).
\]
Multiplying gives the literal word $P=aBfde$. Since $B=CbC$ and
$f=Czk^{-1}$, the selected cone gives
\[
 cP=abzk^{-1}de=abzde k^{-1}.
\]
The last equality uses the actual relations $[k,d]=[k,e]=1$.
The complete regional presentation also has
$[k,a]=[k,b]=[k,z]=1$. Consequently
\begin{equation}\label{eq:B118-boundary-commutes}
 [k,T]=1,\qquad c^2=1.
\end{equation}
These equations give a presentation map from the abstract group
$C_2(c)*\mathbb Z^2(k,T)$ onto the actual subgroup $H$.

Kill the normal closure of $a,b,d$ in the actual regional group.
The complete generic and node relations of B.112, equivalently the
seventeen-relator presentation in B.115, give exactly the quotient
\begin{equation}\label{eq:B118-retaining-quotient}
 \pi_e:R\twoheadrightarrow R_e
 =\langle C,k,z,e\mid C^2=1,\ [C,z]=[k,z]=[k,e]=1,
                                  \ [e,CeC]=1\rangle.
\end{equation}
All other node relators disappear under this specified filling;
the displayed $e$-node remains. In particular this is a quotient of
the whole region, including the retained base action. In $R_e$ put
\begin{equation}\label{eq:B118-quotient-marking}
 s=ze,\qquad f_e=Czk^{-1},\qquad p=f_e e=Csk^{-1}.
\end{equation}
The last equality uses $[k,e]=1$. The actual marked images are
\begin{equation}\label{eq:B118-actual-images}
 \pi_e(c)=C,\quad \pi_e(k)=k,\quad
 \pi_e(T)=sk^{-1},\quad \pi_e(P)=p,\quad
 \pi_e(K)=\langle f_e,e\rangle.
\end{equation}
We next prove that $H_e=\langle C,k,s\rangle$ is the expected
embedded free product, and establish its full topology directly.

\emph{B.118.2. A single exact HNN comparison.}
Start with the abstract group
\begin{equation}\label{eq:B118-abstract-base}
 H_* =\langle C_0,k,s\mid C_0^2=1,\ [k,s]=1\rangle
      =C_2(C_0)*\mathbb Z^2(k,s).
\end{equation}
Set $C_1=s^{-1}C_0s$ and
\[
 D_0=\langle k,C_0,C_1\rangle\leq H_*.
\]
This subgroup is exactly
\begin{equation}\label{eq:B118-D0}
 D_0=\mathbb Z(k)*C_2(C_0)*C_2(C_1).
\end{equation}
Indeed, rewrite the kernel of the homomorphism $H_*\to\mathbb Z$
sending $s$ to $1$ and $C_0,k$ to $0$, using the transversal
$\{s^n:n\in\mathbb Z\}$. The conjugates of $k$ all agree because
$[k,s]=1$. The remaining generators are
$C_n=s^{-n}C_0s^n$, with precisely the relations $C_n^2=1$.
Thus that kernel is the free product of $\mathbb Z(k)$ and all
$C_2(C_n)$, and \eqref{eq:B118-D0} is its indicated full free factor.
In particular both $\langle k,C_1\rangle$ and
$\langle k,C_0\rangle$ embed in $D_0$.

Define the automorphism $\beta$ of $D_0$ by
\begin{equation}\label{eq:B118-beta}
 \beta(k)=k,\qquad
 \beta(C_0)=C_0C_1C_0,\qquad \beta(C_1)=C_0.
\end{equation}
Its inverse fixes $k$ and sends $C_0$ to $C_1$ and $C_1$ to
$C_1C_0C_1$. Both compositions are the identity. With
$r=C_1C_0$, one also has $\beta(r)=r$.
The base itself has the HNN presentation
\begin{equation}\label{eq:B118-base-HNN}
 H_*=
 \langle D_0,s\mid sC_1s^{-1}=C_0,\ [s,k]=1\rangle;
\end{equation}
eliminating $C_1$ gives \eqref{eq:B118-abstract-base}.

There is an exact presentation isomorphism
\begin{equation}\label{eq:B118-retaining-HNN}
 R_e\cong
 \langle H_*,e\mid e x e^{-1}=\beta(x)\ (x\in D_0)\rangle
 \cong H_* *_{D_0}\bigl(D_0\rtimes_{\beta}\mathbb Z(e)\bigr).
\end{equation}
Here the actual $C$ is $C_0$, and the actual $s$ is $ze$.
To verify every substitution, put $z=se^{-1}$.
The relations $[k,z]=[k,e]=1$ are equivalent to
$[k,s]=[k,e]=1$. The relation $[C_0,z]=1$ is
\[
 e^{-1}C_0e=C_1,\qquad\text{equivalently }eC_1e^{-1}=C_0.
\]
Using $C_0e=eC_1$, the equation $[e,C_0eC_0]=1$ becomes
$e^2r=ere$, hence $[e,r]=1$. Together these give
\[
 eC_0e^{-1}=C_0r=C_0C_1C_0,
 \qquad eC_1e^{-1}=C_0,
\]
which are exactly \eqref{eq:B118-beta}; conversely they recover
the original two equations. All substitutions are reversible.
Since $D_0$ is embedded in $H_*$ and $\beta$ is an automorphism,
the ordinary HNN normal-form theorem embeds $H_*$ in
\eqref{eq:B118-retaining-HNN}. Its image is the actual $H_e$.

The composite of the presentation map following
\eqref{eq:B118-boundary-commutes} with $\pi_e$ sends
$c,k,T$ to $C_0,k,sk^{-1}$. This is an isomorphism onto $H_e$:
$s=(sk^{-1})k$ is the inverse change of basis in the abelian factor.
The original presentation map is therefore injective. This proves
\eqref{eq:B118-H}, and $\pi_e|_H:H\to H_e$ is an isomorphism.

\emph{B.118.3. Full induced topology of the HNN base.}
We give a finite-target argument for $H_e\leq R_e$ rather than
assuming a virtual retraction. Let $\psi:H_*\to F$ be an arbitrary
homomorphism to a finite group. The finitely generated group $D_0$
has a characteristic finite-index subgroup
\[
 N\leq\ker(\psi|_{D_0}).
\]
For example, intersect all subgroups of the same index as this kernel;
there are finitely many because $D_0$ is finitely generated.
Thus $\beta(N)=N$. Write $Q=D_0/N$, let $\bar\beta$ be the
induced automorphism, and let $A_1,A_0\leq Q$ be the images of
$\langle k,C_1\rangle,\langle k,C_0\rangle$. Then
$\bar\beta:A_1\to A_0$ is an isomorphism. The quotient
\begin{equation}\label{eq:B118-finite-base-HNN}
 H_N=\langle Q,s\mid s x s^{-1}=\bar\beta(x)\ (x\in A_1)\rangle
\end{equation}
is a proper ordinary HNN extension, with $Q$ embedded, and $\psi$
factors through it. If $m$ is the finite order of $\bar\beta$,
\eqref{eq:B118-retaining-HNN} maps onto
\begin{equation}\label{eq:B118-finite-vertex-amalgam}
 H_N *_{Q}\bigl(Q\rtimes_{\bar\beta}C_m\bigr).
\end{equation}

Only two elementary finite permutation constructions are needed.
First represent $Q$ by its regular action on a finite set. Its
restrictions to $A_1$ and $A_0$ are free actions with the same number
of orbits under their specified isomorphism. An intertwining
permutation realizes the stable letter in
\eqref{eq:B118-finite-base-HNN}. This gives a finite permutation
image of $H_N$ which is injective on $Q$. Combine that image with
the factorization of $\psi$ and denote the resulting finite image
by $F'$. Then $Q\hookrightarrow F'$, and $F'$ retains $\psi$.

Second, for any two finite groups with a common embedded subgroup,
choose free actions on the same finite number of points, divisible
by both group orders. The restrictions to the common subgroup are
unions of regular actions of equal cardinality, so conjugating one
action aligns them. The resulting permutation representation of
their amalgam is injective on both finite vertex groups. Apply this
to
\[
 F' *_{Q}\bigl(Q\rtimes_{\bar\beta}C_m\bigr).
\]
Composing its finite permutation image with
\eqref{eq:B118-finite-vertex-amalgam} gives a finite quotient of
$R_e$ whose kernel on $H_*$ is contained in $\ker\psi$.
This proves full induced profinite topology of $H_e$ in $R_e$.
No residual-finiteness theorem for an infinite amalgam is required.

Pull back these finite quotients along the actual $\pi_e$ and use
the isomorphism $\pi_e|_H$. Every finite quotient of the actual $H$
is then refined by a finite quotient of $R$. Hence $H$ has its full
induced topology in $R$. In particular the completed map
\begin{equation}\label{eq:B118-completed-base-map}
 \widehat\pi_e|_{\widehat H}:
       \widehat H\xrightarrow{\ \sim\ }\widehat{H_e}
          \hookrightarrow\widehat{R_e}
\end{equation}
is injective. This is a comparison of explicitly embedded full
completions, not an identification of an unspecified kernel completion.

\emph{B.118.4. A genuine quotient retraction and both intersections.}
Let $F(x,y)$ be an abstract free group. Define
\begin{equation}\label{eq:B118-quotient-retraction}
 \begin{aligned}
 \tau:R_e&\longrightarrow F(x,y),
   &C,k&\longmapsto1,&z&\longmapsto x,&e&\longmapsto y,\\
 \sigma:F(x,y)&\longrightarrow R_e,
   &x&\longmapsto f_e=Czk^{-1},&y&\longmapsto e.
 \end{aligned}
\end{equation}
Every defining relator of $R_e$ is respected by $\tau$; in
particular $[e,CeC]$ maps to $[y,y]=1$. The map $\sigma$ exists
because its domain is free, and $\tau\sigma$ is the identity.
Thus $K_e=\langle f_e,e\rangle$ is an embedded free group of rank
two, and $\rho=\sigma\tau:R_e\to K_e$ is an actual retraction.
It fixes $f_e,e$, kills $C,k$, and sends $z$ to the literal $f_e$.
Consequently
\begin{equation}\label{eq:B118-retraction-image}
 \rho(H_e)=\langle p\rangle,\qquad p=f_e e=Csk^{-1}.
\end{equation}
Since $p$ belongs to both $K_e$ and $H_e$, applying $\rho$ proves
\[
 K_e\cap H_e=\langle p\rangle.
\]

The same retraction gives $K_e$ its full induced topology in $R_e$.
The element $p$ is a member of the free basis $(p,e)$ of $K_e$,
so its cyclic subgroup also has its full topology. If
$v\in\widehat{K_e}\cap\widehat{H_e}$ inside $\widehat{R_e}$,
then $\widehat\rho(v)=v$. The image
$\widehat\rho(\widehat{H_e})$ is exactly
$\overline{\langle p\rangle}$: it is compact and contains the
dense image of $H_e$, which is $\langle p\rangle$.
The reverse containment is immediate. Therefore
\begin{equation}\label{eq:B118-quotient-completed-intersection}
 \widehat{K_e}\cap\widehat{H_e}
       =\overline{\langle p\rangle}\quad\text{in }\widehat{R_e}.
\end{equation}
This proves the completed intersection directly; it is not deduced
from the discrete equality.

The actual quotient sends $K$ onto $K_e$ and $H$ isomorphically onto
$H_e$, with $P$ mapping to $p$. The discrete intersection in
\eqref{eq:B118-intersections} follows at once. For the completed
one, B.114/B.117 gives the actual embedding
$\widehat K\hookrightarrow\widehat R$; by compactness and density
its image under $\widehat\pi_e$ is all of $\widehat{K_e}$.
If $v\in\widehat K\cap\widehat H$, then
\eqref{eq:B118-quotient-completed-intersection} puts
$\widehat\pi_e(v)$ in $\overline{\langle p\rangle}$.
The compact image of $\overline{\langle P\rangle}$ is exactly
$\overline{\langle p\rangle}$, so choose
$w\in\overline{\langle P\rangle}$ with
$\widehat\pi_e(w)=\widehat\pi_e(v)$. Both $v,w$ lie in
$\widehat H$, where \eqref{eq:B118-completed-base-map} is
injective. Hence $v=w$. The reverse containment follows from
$P\in K\cap H$. Finally, the map $H\to\langle P\rangle$
sending $c,k$ to $1$ and $T$ to the actual $P$ is a retraction;
it fixes $P=cT$. This identifies its compact closure with its full
procyclic completion.

\emph{B.118.5. The corner-side embedding and topology.}
The group in \eqref{eq:B118-corner} is the ordinary amalgam
\[
 \bigl((C_2(D)*\mathbb Z(u))\times\mathbb Z(k)\bigr)
       *_{\mathbb Z(u)}
       \bigl(\mathbb Z(u)\rtimes C_2(j)\bigr),
\]
with the second factor acting on $u$ by inversion. Both cyclic edge
maps are injective. For the auxiliary parity taking $j$ to $1$ and
$D,u,k$ to $0$, the index-two subgroup has the exact presentation
\begin{equation}\label{eq:B118-corner-even}
 W_{\partial}^{+}=
 \langle D,k,u,k',c\mid
 D^2=c^2=1,\ [D,k]=[k,u]=[u,k']=[k',c]=1\rangle,
\end{equation}
where $k'=jkj$ and $c=jDj$. This follows by the two-coset
rewrite: the second $u$ is $juj=u^{-1}$, and the two copies of
each generic commutator are the four displayed edges.
This parity is not being identified with the original ordinary cover.

In \eqref{eq:B118-corner-even}, $T=uD$ commutes with $k$.
There are mutually compatible maps
\[
 C_2(c)*\mathbb Z^2(k,T)\longrightarrow W_{\partial}^{+},
 \qquad
 W_{\partial}^{+}\longrightarrow C_2(c)*\mathbb Z^2(k,T),
\]
the first sending $T$ to $uD$ and fixing $c,k$, and the second
sending $D,k'$ to $1$, $u$ to $T$, and fixing $c,k$.
Every edge relation is respected, and their composite on the free
product is the identity. Thus the first map is injective and its
image is an actual retract. Retractions give full induced profinite
topology by pulling back every finite quotient. Passing through the
index-two inclusion $W_{\partial}^{+}\leq W_{\partial}$, and
taking a finite-index core when necessary, gives the claimed full
topology in the whole corner.

The two full-topology statements and
\eqref{eq:B118-intersections} concern the specified regional and
corner groups. They do not prove properness of a new full profinite
amalgam over $H$, nor injection after the four remaining prefix
mirror relations or the other actual regional relations are imposed.
No source-centralizer or main decision is inferred from this comparison.
\end{proof}

\paragraph{Proof provenance and exact markings.}
The proof uses the actual upper-gap cone, the full boundary word and
the literal quotient killing $a,b,d$ throughout. The earlier pentagon
proof is preserved as an independent derivation; the single HNN
comparison above suffices here. The completed intersection uses the
actual retraction and injective completed quotient restriction, with
no exchange of completion and intersection. The underlying reviewed
records are
\begingroup\raggedright
\path{records/pentagon_hnn_boundary_proof_candidate24.md},
\path{records/completed_boundary_neck_intersection24.md},
\path{work_DE/next_marginal_source/review_pentagon_boundary_full24.md},
\path{work_DE/next_marginal_source/actual_mirror_triangle_and_boundary_corner24.md}.
\par\endgroup
